% GENERATED FILE. Edit canonical lessons, then run npm run generate:books.
% canonical-source-hash: fnv1a64:8991eaa1d36c2ed3
% canonical-lessons: KA-C332-vipattu, KA-C332-siluku, KA-C332-badukuli, KA-C332-santrasta, KA-C332-munsucane

\chapter{After the Disaster}
\label{ch:ka-after-the-disaster}
\hlchaptermodality{332}

\begin{chapteropening}
\textbf{By the end of this chapter:} \emph{I can say a disaster, to be caught, to survive, a victim and a forewarning.}

Complete the last lesson of chapter 332: I can say a disaster, to be caught, to survive, a victim and a forewarning.
\end{chapteropening}

\section[\texorpdfstring{\kn{ವಿಪತ್ತು} (vipattu)}{ವಿಪತ್ತು (vipattu)}]{\kn{ವಿಪತ್ತು} (vipattu) — a disaster; a calamity}

\label{lesson:KA-C332-vipattu}



\begin{quote}
Before the new one: say the Kannada for ancient, then the Kannada for an era.
\end{quote}

\subsection*{You'll want to know: \kn{ವಿಪತ್ತು}}
\textbf{\kn{ವಿಪತ್ತು}} — \emph{vipattu} — \textquotedblleft{}a disaster; a calamity\textquotedblright{}.

\begin{quote}
\kn{ಪ್ರವಾಹ} \kn{ಆ} \kn{ಜಿಲ್ಲೆಗೆ} \kn{ದೊಡ್ಡ} \kn{ವಿಪತ್ತು} \kn{ತಂದಿತು}.
\end{quote}

\begin{quote}
\emph{pravāha ā jillege doḍḍa vipattu tanditu.}
\end{quote}

\begin{quote}
The flood brought a great disaster to that district.
\end{quote}

Plural: \kn{ವಿಪತ್ತುಗಳು}, \emph{vipattugaḷu}.

\kn{ಪ್ರಕೃತಿ} \kn{ವಿಕೋಪ} is the news phrase for a natural disaster.

\begin{cousinweb}
\textbf{\kn{ವಿಪತ್ತು}} (\emph{vipattu}) is borrowed from Sanskrit \emph{vipatti}, \textquotedblleft{}misfortune, adversity, disaster\textquotedblright{} (CDIAL 11788).
\end{cousinweb}

\begin{grammarlens}[title={What you've built}]
One more word for telling what happened.
\end{grammarlens}

\subsection*{Your turn}
\begin{itemize}
\raggedright
  \item \emph{Say it:} \emph{vipattu}
  \item \emph{Say it:} \emph{vipattu}, once more
  \item \emph{Say it:} \emph{pravāha ā jillege doḍḍa vipattu tanditu.} — The flood brought a great disaster to that district.
  \item \emph{Recall:} say the Kannada for ancient, then the Kannada for an era, then say \emph{vipattu} again
\end{itemize}

\begin{tcolorbox}[breakable,colback=teal!4,colframe=teal!35!black,title={Before you move on}]
What does \kn{ವಿಪತ್ತು} mean? (\textquotedblleft{}A disaster; a calamity\textquotedblright{}.) Say it once more.
\end{tcolorbox}
\section[\texorpdfstring{\kn{ಸಿಲುಕು} (siluku)}{ಸಿಲುಕು (siluku)}]{\kn{ಸಿಲುಕು} (siluku) — to be caught; to be trapped}

\label{lesson:KA-C332-siluku}



\begin{quote}
Before the new one: say the Kannada for an era, then the Kannada for a disaster.
\end{quote}

\subsection*{You'll want to know: \kn{ಸಿಲುಕು}}
\textbf{\kn{ಸಿಲುಕು}} — \emph{siluku} — \textquotedblleft{}to be caught; to be trapped\textquotedblright{}.

\begin{quote}
\kn{ಪ್ರವಾಹದಲ್ಲಿ} \kn{ನೂರು} \kn{ಜನ} \kn{ಸಿಲುಕಿದರು}.
\end{quote}

\begin{quote}
\emph{pravāhadalli nūru jana silukidaru.}
\end{quote}

\begin{quote}
A hundred people were trapped in the flood.
\end{quote}

Past (I …): \kn{ಸಿಲುಕಿದೆ}, \emph{silukide}.

Where one is trapped takes -\kn{ಅಲ್ಲಿ}: \kn{ಮನೆಯಲ್ಲಿ} \kn{ಸಿಲುಕು}.

\begin{cousinweb}
\textbf{\kn{ಸಿಲುಕು}} (\emph{siluku}) is native Dravidian: DEDR 2498 lists Kannada \emph{siluku} beside Tamil \emph{cikku}, \textquotedblleft{}to be entangled, caught, ensnared\textquotedblright{}.
\end{cousinweb}

\begin{grammarlens}[title={What you've built}]
One more word for telling what happened.
\end{grammarlens}

\subsection*{Your turn}
\begin{itemize}
\raggedright
  \item \emph{Say it:} \emph{siluku}
  \item \emph{Say it:} \emph{siluku}, once more
  \item \emph{Say it:} \emph{pravāhadalli nūru jana silukidaru.} — A hundred people were trapped in the flood.
  \item \emph{Recall:} say the Kannada for an era, then the Kannada for a disaster, then say \emph{siluku} again
\end{itemize}

\begin{tcolorbox}[breakable,colback=teal!4,colframe=teal!35!black,title={Before you move on}]
What does \kn{ಸಿಲುಕು} mean? (\textquotedblleft{}To be caught; to be trapped\textquotedblright{}.) Say it once more.
\end{tcolorbox}
\section[\texorpdfstring{\kn{ಬದುಕುಳಿ} (badukuḷi)}{ಬದುಕುಳಿ (badukuḷi)}]{\kn{ಬದುಕುಳಿ} (badukuḷi) — to survive}

\label{lesson:KA-C332-badukuli}



\begin{quote}
Before the new one: say the Kannada for a disaster, then the Kannada for to be caught.
\end{quote}

\subsection*{You'll want to know: \kn{ಬದುಕುಳಿ}}
\textbf{\kn{ಬದುಕುಳಿ}} — \emph{badukuḷi} — \textquotedblleft{}to survive\textquotedblright{}.

\begin{quote}
\kn{ಆ} \kn{ಅಪಘಾತದಲ್ಲಿ} \kn{ಒಬ್ಬನೇ} \kn{ಬದುಕುಳಿದನು}.
\end{quote}

\begin{quote}
\emph{ā apaghātadalli obbanē badukuḷidanu.}
\end{quote}

\begin{quote}
Only one man survived that accident.
\end{quote}

Past (I …): \kn{ಬದುಕುಳಿದೆ}, \emph{badukuḷide}.

To be left alive after danger; \kn{ಉಳಿ} is to remain.

\begin{cousinweb}
\textbf{\kn{ಬದುಕುಳಿ}} (\emph{badukuḷi}) is a native Dravidian compound: \emph{baduku}, \textquotedblleft{}to live\textquotedblright{} (DEDR 5372), and \emph{uḷi}, \textquotedblleft{}to remain\textquotedblright{} (DEDR 1009).
\end{cousinweb}

\begin{grammarlens}[title={What you've built}]
One more word for telling what happened.
\end{grammarlens}

\subsection*{Your turn}
\begin{itemize}
\raggedright
  \item \emph{Say it:} \emph{badukuḷi}
  \item \emph{Say it:} \emph{badukuḷi}, once more
  \item \emph{Say it:} \emph{ā apaghātadalli obbanē badukuḷidanu.} — Only one man survived that accident.
  \item \emph{Recall:} say the Kannada for a disaster, then the Kannada for to be caught, then say \emph{badukuḷi} again
\end{itemize}

\begin{tcolorbox}[breakable,colback=teal!4,colframe=teal!35!black,title={Before you move on}]
What does \kn{ಬದುಕುಳಿ} mean? (\textquotedblleft{}To survive\textquotedblright{}.) Say it once more.
\end{tcolorbox}
\section[\texorpdfstring{\kn{ಸಂತ್ರಸ್ತ} (santrasta)}{ಸಂತ್ರಸ್ತ (santrasta)}]{\kn{ಸಂತ್ರಸ್ತ} (santrasta) — a victim; one affected}

\label{lesson:KA-C332-santrasta}



\begin{quote}
Before the new one: say the Kannada for to be caught, then the Kannada for to survive.
\end{quote}

\subsection*{You'll want to know: \kn{ಸಂತ್ರಸ್ತ}}
\textbf{\kn{ಸಂತ್ರಸ್ತ}} — \emph{santrasta} — \textquotedblleft{}a victim; one affected\textquotedblright{}.

\begin{quote}
\kn{ಸಂತ್ರಸ್ತರಿಗೆ} \kn{ಊಟ} \kn{ಮತ್ತು} \kn{ಬಟ್ಟೆ} \kn{ಕೊಟ್ಟರು}.
\end{quote}

\begin{quote}
\emph{santrastarige ūṭa mattu baṭṭe koṭṭaru.}
\end{quote}

\begin{quote}
They gave food and clothes to the victims.
\end{quote}

Plural: \kn{ಸಂತ್ರಸ್ತರು}, \emph{santrastaru}.

A news word for people hit by a flood, a fire or a quake: \kn{ಪ್ರವಾಹ} \kn{ಸಂತ್ರಸ್ತರು}.

\begin{cousinweb}
\textbf{\kn{ಸಂತ್ರಸ್ತ}} (\emph{santrasta}) is borrowed from Sanskrit \emph{saṃtrasta}, \textquotedblleft{}trembling with fear, alarmed\textquotedblright{} (Monier-Williams).
\end{cousinweb}

\begin{grammarlens}[title={What you've built}]
One more word for telling what happened.
\end{grammarlens}

\subsection*{Your turn}
\begin{itemize}
\raggedright
  \item \emph{Say it:} \emph{santrasta}
  \item \emph{Say it:} \emph{santrasta}, once more
  \item \emph{Say it:} \emph{santrastarige ūṭa mattu baṭṭe koṭṭaru.} — They gave food and clothes to the victims.
  \item \emph{Recall:} say the Kannada for to be caught, then the Kannada for to survive, then say \emph{santrasta} again
\end{itemize}

\begin{tcolorbox}[breakable,colback=teal!4,colframe=teal!35!black,title={Before you move on}]
What does \kn{ಸಂತ್ರಸ್ತ} mean? (\textquotedblleft{}A victim; one affected\textquotedblright{}.) Say it once more.
\end{tcolorbox}
\section[\texorpdfstring{\kn{ಮುನ್ಸೂಚನೆ} (munsūcane)}{ಮುನ್ಸೂಚನೆ (munsūcane)}]{\kn{ಮುನ್ಸೂಚನೆ} (munsūcane) — a forewarning; a forecast}

\label{lesson:KA-C332-munsucane}



\begin{quote}
Before the new one: say the Kannada for to survive, then the Kannada for a victim.
\end{quote}

\subsection*{You'll want to know: \kn{ಮುನ್ಸೂಚನೆ}}
\textbf{\kn{ಮುನ್ಸೂಚನೆ}} — \emph{munsūcane} — \textquotedblleft{}a forewarning; a forecast\textquotedblright{}.

\begin{quote}
\kn{ಹವಾಮಾನ} \kn{ಇಲಾಖೆ} \kn{ಮೂರು} \kn{ದಿನದ} \kn{ಮುಂಚೆ} \kn{ಮುನ್ಸೂಚನೆ} \kn{ಕೊಟ್ಟಿತ್ತು}.
\end{quote}

\begin{quote}
\emph{havāmāna ilākhe mūru dinada muñce munsūcane koṭṭittu.}
\end{quote}

\begin{quote}
The weather department had given a warning three days earlier.
\end{quote}

\kn{ಮುನ್}, before, and \kn{ಸೂಚನೆ}, a notice.

\begin{cousinweb}
\textbf{\kn{ಮುನ್ಸೂಚನೆ}} (\emph{munsūcane}) is a Kannada compound of native \emph{mun}, \textquotedblleft{}before, in front\textquotedblright{} (DEDR 5020), and \kn{ಸೂಚನೆ}, from Sanskrit \emph{sūcanā}, \textquotedblleft{}an indication, a communication\textquotedblright{} (Monier-Williams).
\end{cousinweb}

\begin{grammarlens}[title={What you've built}]
One more word for telling what happened.
\end{grammarlens}

\subsection*{Your turn}
\begin{itemize}
\raggedright
  \item \emph{Say it:} \emph{munsūcane}
  \item \emph{Say it:} \emph{munsūcane}, once more
  \item \emph{Say it:} \emph{havāmāna ilākhe mūru dinada muñce munsūcane koṭṭittu.} — The weather department had given a warning three days earlier.
  \item \emph{Recall:} say the Kannada for to survive, then the Kannada for a victim, then say \emph{munsūcane} again
\end{itemize}

\begin{tcolorbox}[breakable,colback=teal!4,colframe=teal!35!black,title={Before you move on}]
What does \kn{ಮುನ್ಸೂಚನೆ} mean? (\textquotedblleft{}A forewarning; a forecast\textquotedblright{}.) Say it once more.
\end{tcolorbox}
